\documentclass{6160}

\author{Guanyuming He}
% \problem{A} means Problem Set A.
\problem{0}


\collab{none}
% or give names, e.g., \collab{Alyssa P. Hacker and A. Student}

\begin{document}

\section*{Problem 1}

A) In normal logging in programs, there are attempt rate limits, which would
become useless should the hashes are not protected.

B) I suppose there is no real gain in limiting which hash function to use for
all users' passwords, except that it will make the system design simpler.  In
contrast, if the hash functions are unlimited, the system can be easily upgraded
in the future with safer hash algorithms (without the user's even knowing about
it!)

C) Making a hash evaluation much slower, thus limiting brute force attacks.

D) I would set a reasonable count that can make it harder at least 1000 times
for an attacker, while not enough to cause big slow down to the server (ideally,
maybe 100th of a second for one authentication), since the server is popular.
If 100 people are using the server together, that will take a total of 1 second
of CPU time, which is big already.

E) I will use a different strategy on my personal laptop, because only me will
be using it, and I don't care if I will wait for a few seconds for one mistyped
password, I will set the count so that it will take at least seconds to
calcaulate such a hash.

% make sure you have each problem in a different page.
\newpage \section*{Problem 2}

B)

D)

% make sure you have each problem in a different page.
\newpage \section*{Problem 4}

A)

\end{document}


